% GENERATED FILE. Edit canonical lessons, then run npm run generate:books.
% canonical-source-hash: fnv1a64:e7257021edc81743
% canonical-lessons: HI-C469-krpan, HI-C469-mitvyayi, HI-C469-dan, HI-C469-lobhi, HI-C469-bhikhari

\chapter{Generous and Mean}
\label{ch:hi-generous-and-mean}
\hlchaptermodality{469}

\begin{chapteropening}
\textbf{By the end of this chapter:} \emph{I can say whether someone is generous or mean.}

Complete the last lesson of chapter 469: I can say whether someone is generous or mean.
\end{chapteropening}

\section[\texorpdfstring{\hi{कृपण} (kṛpaṇ)}{कृपण (kṛpaṇ)}]{\hi{कृपण} (kṛpaṇ) — miserly; tight-fisted}

\label{lesson:HI-C469-krpan}



\begin{quote}
Before the new one: say the Hindi for to snap at, then the Hindi for to snap irritably.
\end{quote}

\subsection*{You'll want to know: \hi{कृपण}}
\textbf{\hi{कृपण}} — \emph{kṛpaṇ} — \textquotedblleft{}miserly; tight-fisted\textquotedblright{}.

\begin{quote}
\hi{मेरा} \hi{पड़ोसी} \hi{बहुत} \hi{कृपण} \hi{है}, \hi{अपने} \hi{बच्चों} \hi{को} \hi{भी} \hi{पैसे} \hi{नहीं} \hi{देता।}
\end{quote}

\begin{quote}
\emph{merā paṛosī bahut kṛpaṇ hai, apne baccoṁ ko bhī paise nahīṃ detā.}
\end{quote}

\begin{quote}
My neighbour is terribly tight-fisted; he won't give money even to his own children.
\end{quote}

\begin{cousinweb}
\textbf{\hi{कृपण}} (\emph{kṛpaṇ}) is a tatsama, borrowed from Sanskrit \emph{kṛpaṇa}, first \textquotedblleft{}pitiable\textquotedblright{} and then \textquotedblleft{}miserly, stingy\textquotedblright{} (Monier-Williams). It is a formal word, the one books use for a miser.
\end{cousinweb}

\begin{grammarlens}[title={What you've built}]
One more word for describing what you lived and felt.
\end{grammarlens}

\subsection*{Your turn}
\begin{itemize}
\raggedright
  \item \emph{Say it:} \emph{kṛpaṇ}
  \item \emph{Say it:} \emph{kṛpaṇ}, once more
  \item \emph{Say it:} \emph{merā paṛosī bahut kṛpaṇ hai, apne baccoṁ ko bhī paise nahīṃ detā.} — My neighbour is terribly tight-fisted; he won't give money even to his own children.
  \item \emph{Recall:} say the Hindi for to snap at, then the Hindi for to snap irritably, then say \emph{kṛpaṇ} again
\end{itemize}

\begin{tcolorbox}[breakable,colback=teal!4,colframe=teal!35!black,title={Before you move on}]
What does \hi{कृपण} mean? (\textquotedblleft{}Miserly; tight-fisted\textquotedblright{}.) Say it once more.
\end{tcolorbox}
\section[\texorpdfstring{\hi{मितव्ययी} (mitvyayī)}{मितव्ययी (mitvyayī)}]{\hi{मितव्ययी} (mitvyayī) — thrifty; careful with money}

\label{lesson:HI-C469-mitvyayi}



\begin{quote}
Before the new one: say the Hindi for to snap irritably, then the Hindi for miserly.
\end{quote}

\subsection*{You'll want to know: \hi{मितव्ययी}}
\textbf{\hi{मितव्ययी}} — \emph{mitvyayī} — \textquotedblleft{}thrifty; careful with money\textquotedblright{}.

\begin{quote}
\hi{दादी} \hi{बहुत} \hi{मितव्ययी} \hi{हैं}, \hi{एक} \hi{पैसा} \hi{भी} \hi{बेकार} \hi{नहीं} \hi{करतीं।}
\end{quote}

\begin{quote}
\emph{dādī bahut mitvyayī haiṁ, ek paisā bhī bekār nahīṃ kartīṃ.}
\end{quote}

\begin{quote}
Grandmother is very thrifty; she doesn't waste a single paisa.
\end{quote}

\begin{cousinweb}
\textbf{\hi{मितव्ययी}} (\emph{mitvyayī}) is a tatsama, borrowed from Sanskrit \emph{mitavyayin}, \textquotedblleft{}spending little, frugal\textquotedblright{}, built on \emph{mita}, \textquotedblleft{}measured\textquotedblright{}, and \emph{vyaya}, \textquotedblleft{}spending\textquotedblright{} (Monier-Williams). \emph{kṛpaṇ} is mean; \emph{mitvyayī} is sensible. It is a formal word; speech often says someone is careful with money.
\end{cousinweb}

\begin{grammarlens}[title={What you've built}]
One more word for describing what you lived and felt.
\end{grammarlens}

\subsection*{Your turn}
\begin{itemize}
\raggedright
  \item \emph{Say it:} \emph{mitvyayī}
  \item \emph{Say it:} \emph{mitvyayī}, once more
  \item \emph{Say it:} \emph{dādī bahut mitvyayī haiṁ, ek paisā bhī bekār nahīṃ kartīṃ.} — Grandmother is very thrifty; she doesn't waste a single paisa.
  \item \emph{Recall:} say the Hindi for to snap irritably, then the Hindi for miserly, then say \emph{mitvyayī} again
\end{itemize}

\begin{tcolorbox}[breakable,colback=teal!4,colframe=teal!35!black,title={Before you move on}]
What does \hi{मितव्ययी} mean? (\textquotedblleft{}Thrifty; careful with money\textquotedblright{}.) Say it once more.
\end{tcolorbox}
\section[\texorpdfstring{\hi{दान} (dān)}{दान (dān)}]{\hi{दान} (dān) — a donation; charity}

\label{lesson:HI-C469-dan}



\begin{quote}
Before the new one: say the Hindi for miserly, then the Hindi for thrifty.
\end{quote}

\subsection*{You'll want to know: \hi{दान}}
\textbf{\hi{दान}} — \emph{dān} — \textquotedblleft{}a donation; charity\textquotedblright{}.

\begin{quote}
\hi{दादी} \hi{हर} \hi{साल} \hi{मंदिर} \hi{में} \hi{दान} \hi{देती} \hi{हैं।}
\end{quote}

\begin{quote}
\emph{dādī har sāl mandir meṁ dān detī haiṁ.}
\end{quote}

\begin{quote}
Grandmother gives to the temple every year.
\end{quote}

Masculine.

\begin{cousinweb}
\textbf{\hi{दान}} (\emph{dān}) is a tatsama, borrowed from Sanskrit \emph{dāna}, \textquotedblleft{}the act of giving\textquotedblright{}, from \emph{dā}, \textquotedblleft{}to give\textquotedblright{} (Monier-Williams), the same root as \emph{denā}.
\end{cousinweb}

\begin{grammarlens}[title={What you've built}]
One more word for describing what you lived and felt.
\end{grammarlens}

\subsection*{Your turn}
\begin{itemize}
\raggedright
  \item \emph{Say it:} \emph{dān}
  \item \emph{Say it:} \emph{dān}, once more
  \item \emph{Say it:} \emph{dādī har sāl mandir meṁ dān detī haiṁ.} — Grandmother gives to the temple every year.
  \item \emph{Recall:} say the Hindi for miserly, then the Hindi for thrifty, then say \emph{dān} again
\end{itemize}

\begin{tcolorbox}[breakable,colback=teal!4,colframe=teal!35!black,title={Before you move on}]
What does \hi{दान} mean? (\textquotedblleft{}A donation; charity\textquotedblright{}.) Say it once more.
\end{tcolorbox}
\section[\texorpdfstring{\hi{लोभी} (lobhī)}{लोभी (lobhī)}]{\hi{लोभी} (lobhī) — greedy; grasping}

\label{lesson:HI-C469-lobhi}



\begin{quote}
Before the new one: say the Hindi for thrifty, then the Hindi for a donation.
\end{quote}

\subsection*{You'll want to know: \hi{लोभी}}
\textbf{\hi{लोभी}} — \emph{lobhī} — \textquotedblleft{}greedy; grasping\textquotedblright{}.

\begin{quote}
\hi{लोभी} \hi{आदमी} \hi{को} \hi{कभी} \hi{संतोष} \hi{नहीं} \hi{होता।}
\end{quote}

\begin{quote}
\emph{lobhī ādmī ko kabhī santoṣ nahīṃ hotā.}
\end{quote}

\begin{quote}
A greedy man is never content.
\end{quote}

\begin{cousinweb}
\textbf{\hi{लोभी}} (\emph{lobhī}) is a tatsama, borrowed from Sanskrit \emph{lobhin}, \textquotedblleft{}covetous, avaricious\textquotedblright{} (Monier-Williams), from \emph{lobha}, \textquotedblleft{}greed\textquotedblright{}.
\end{cousinweb}

\begin{grammarlens}[title={What you've built}]
One more word for describing what you lived and felt.
\end{grammarlens}

\subsection*{Your turn}
\begin{itemize}
\raggedright
  \item \emph{Say it:} \emph{lobhī}
  \item \emph{Say it:} \emph{lobhī}, once more
  \item \emph{Say it:} \emph{lobhī ādmī ko kabhī santoṣ nahīṃ hotā.} — A greedy man is never content.
  \item \emph{Recall:} say the Hindi for thrifty, then the Hindi for a donation, then say \emph{lobhī} again
\end{itemize}

\begin{tcolorbox}[breakable,colback=teal!4,colframe=teal!35!black,title={Before you move on}]
What does \hi{लोभी} mean? (\textquotedblleft{}Greedy; grasping\textquotedblright{}.) Say it once more.
\end{tcolorbox}
\section[\texorpdfstring{\hi{भिखारी} (bhikhārī)}{भिखारी (bhikhārī)}]{\hi{भिखारी} (bhikhārī) — a beggar}

\label{lesson:HI-C469-bhikhari}



\begin{quote}
Before the new one: say the Hindi for a donation, then the Hindi for greedy.
\end{quote}

\subsection*{You'll want to know: \hi{भिखारी}}
\textbf{\hi{भिखारी}} — \emph{bhikhārī} — \textquotedblleft{}a beggar\textquotedblright{}.

\begin{quote}
\hi{मंदिर} \hi{के} \hi{सामने} \hi{एक} \hi{भिखारी} \hi{बैठा} \hi{था।}
\end{quote}

\begin{quote}
\emph{mandir ke sāmne ek bhikhārī baiṭhā thā.}
\end{quote}

\begin{quote}
A beggar was sitting in front of the temple.
\end{quote}

Masculine.

\begin{cousinweb}
\textbf{\hi{भिखारी}} (\emph{bhikhārī}) is a tadbhava: CDIAL lists Hindi \emph{bhikhārī}, \textquotedblleft{}mendicant\textquotedblright{}, under Sanskrit \emph{bhikṣācara}, \textquotedblleft{}one who goes about begging\textquotedblright{} (CDIAL 9486). The feminine is \emph{bhikhārin}.
\end{cousinweb}

\begin{grammarlens}[title={What you've built}]
One more word for describing what you lived and felt.
\end{grammarlens}

\subsection*{Your turn}
\begin{itemize}
\raggedright
  \item \emph{Say it:} \emph{bhikhārī}
  \item \emph{Say it:} \emph{bhikhārī}, once more
  \item \emph{Say it:} \emph{mandir ke sāmne ek bhikhārī baiṭhā thā.} — A beggar was sitting in front of the temple.
  \item \emph{Recall:} say the Hindi for a donation, then the Hindi for greedy, then say \emph{bhikhārī} again
\end{itemize}

\begin{tcolorbox}[breakable,colback=teal!4,colframe=teal!35!black,title={Before you move on}]
What does \hi{भिखारी} mean? (\textquotedblleft{}A beggar\textquotedblright{}.) Say it once more.
\end{tcolorbox}
